%! Modeling with Exponential Functions: Interest, Half-Life, and e — Unit 6, Lesson 6.4 — Warm-Up (Answer Key)
\documentclass[10pt]{article}
\usepackage{algebra2-article}
\usepackage{algebra2-key}   % pulls in -boxes; do NOT also load -boxes
\newcommand{\TallMath}[1]{$\displaystyle #1\rule[-1.4em]{0pt}{3.2em}$}

\begin{document}

\pageheader{Unit 6, Lesson 6.4}{Warm-Up --- Answer Key}
\namedateperiod

\vspace{0.10in}

\begin{notesbox}{Getting Ready: Skills We'll Use Today}
\small Today's exponential models come from money and medicine. Nothing below is new: item 1 is Lesson
6.1's percent-to-factor translation, item 2 is Unit 5 exponent work, and item 3 is yesterday.

\vspace{5pt}
\textbf{1. One yearly rate, cut into smaller pieces.} An account pays $6\%$ per year. If interest is
added more than once a year, each payment uses only \emph{part} of that rate.

\vspace{3pt}
\begin{center}
\renewcommand{\arraystretch}{1.4}
\begin{tabular}{|l|c|c|c|c|}
\hline
\textbf{Interest is added} & \textbf{times per year} & \textbf{rate per period} & \textbf{factor per period} & \textbf{periods in 5 years} \\
\hline
once a year & $1$ & $0.06$ & $1.06$ & $5$ \\
\hline
every quarter & $4$ & \ans{$0.015$} & \ans{$1.015$} & \ans{$20$} \\
\hline
every month & $12$ & \ans{$0.005$} & \ans{$1.005$} & \ans{$60$} \\
\hline
\end{tabular}
\end{center}

\vspace{0.20in}

\textbf{2. Halving, over and over.} Evaluate. Round decimals to three places.

\vspace{4pt}
\hspace*{1.2em} $\left(\dfrac12\right)^{3} =$ \ans{$\frac18$} \qquad
$\left(\dfrac12\right)^{0} =$ \ans{$1$} \qquad
$\left(\dfrac12\right)^{1/2} \approx$ \ans{$0.707$} \qquad
$\left(\dfrac12\right)^{1.5} \approx$ \ans{$0.354$}

\vspace{6pt}
\hspace*{1.2em} A medicine loses half of whatever is left every $6$ hours. How many \emph{halvings}
happen in $18$ hours? \ans{$3$}

\vspace{4pt}
\hspace*{1.2em} How many halvings happen in $9$ hours? \ans{$1.5$} \; Explain how a number of
halvings can fail to be a whole number. \ansline{At $9$ hours the sample is partway through its second halving --- $9\div6=1.5$, and a fractional exponent is a perfectly good exponent (Lesson 5.1).}

\vspace{0.22in}

\textbf{3. From yesterday.} Solve using a common base (Lesson 6.3).

\vspace{4pt}
\hspace*{1.2em} $2^{\,t} = 32$ \hfill $t=$ \ans{$5$}

\vspace{6pt}
\hspace*{1.2em} Now look at $\;(1.05)^{\,t} = 2$. \; What would Step 1 ask you to do, and why can you
not do it? \ansline{Step 1 says write both sides as powers of one base --- but $2$ is not a power of $1.05$, so there is no common base to find.}

\end{notesbox}

\begin{teachernote}
\textbf{Item 1 is the whole lesson in miniature, and it should be debriefed at the board.} Two numbers
change together and students routinely change only one: the rate is \emph{divided} by $n$ and the count
of periods is \emph{multiplied} by $n$. Say the sentence the lesson is built on while the table is
still on the board --- \emph{the base is what happens in one period, and the exponent counts the
periods} --- and leave the quarterly row ($0.015$, $1.015$, $20$) written up all period. The two
signature errors of the day are exactly the two halves of that row done separately: $1.06$ kept as the
base with $20$ as the exponent, or $1.015$ used with $5$.

\textbf{Item 2's last two answers are why half-life needs $t/h$ and not $t$.} Students accept $3$
halvings in $18$ hours instantly; $1.5$ halvings in $9$ hours is the one that makes them stop. Do not
resolve it with algebra --- point at $\left(\frac12\right)^{1.5}\approx0.354$, which they just
computed, and note that the answer sits sensibly between $\frac12$ and $\frac14$. That is Lesson 5.1's
rational exponents arriving in a context where they mean something.

\textbf{Item 3 is the wall, arriving before the lesson does.} They will meet
$(1.05)^{t}=2$ again in Section 4 as the closing question, so accept a short answer now and do not
solve it. If a student says ``no solution,'' correct it on the spot --- that is exactly the
misconception yesterday's exit ticket was written to catch, and today's Section 4 depends on the
distinction holding.
\end{teachernote}

\end{document}
